% Telos shared preamble.
%
% Every Telos publication is designed to be printed on a home laser printer,
% one sheet at a time, and carried into a boat or a garage. Two constraints
% follow from that and govern everything below:
%
%   1. Pure monochrome. No value is ever a hue. Emphasis is carried by rule
%      weight, fill, and type, never by color, so a grayscale or black-only
%      printer loses nothing.
%   2. Card density. A "one-pager" is a hard contract: the leaf must fit on a
%      single side of US Letter. The layout is therefore tight by default and
%      the environments below are sized to leave room for a plate.

\documentclass[10pt,letterpaper]{article}

\usepackage[margin=0.55in,includefoot,footskip=0.24in]{geometry}
\usepackage[T1]{fontenc}
\usepackage[utf8]{inputenc}
\usepackage{lmodern}
\usepackage{microtype}
\usepackage{array}
\usepackage{booktabs}
\usepackage{longtable}
\usepackage{tabularx}
\usepackage{enumitem}
\usepackage{needspace}
\usepackage{multicol}
\usepackage{ragged2e}
\usepackage{graphicx}
\usepackage{wrapfig}
\usepackage{xcolor}
\usepackage{fancyhdr}
\usepackage{titlesec}
\usepackage{hyperref}
\usepackage[most]{tcolorbox}
\usepackage{tikz}
\usetikzlibrary{arrows.meta,calc,positioning,shapes.geometric,decorations.pathmorphing,patterns,fadings}

% Semantic names are kept so a document reads intelligibly, but every value is
% black, white, or a neutral gray that survives a black-only print driver.
\definecolor{telosink}{gray}{0}
\definecolor{telospaper}{gray}{1}
\definecolor{telosrule}{gray}{0.35}
\definecolor{telosfill}{gray}{0.90}
\definecolor{telosdeep}{gray}{0.72}
\definecolor{telosmute}{gray}{0.45}

\hypersetup{hidelinks,pdfborder={0 0 0}}

% Keep an unchanged source byte-reproducible across rebuilds: no build-clock
% creation date and no random trailer ID.
\pdfinfoomitdate=1
\pdftrailerid{}

\setlength{\parindent}{0pt}
\setlength{\parskip}{0.42em}
\setlength{\tabcolsep}{4pt}
\setlength{\columnsep}{1.1em}
\renewcommand{\arraystretch}{1.18}
\setlist[itemize]{leftmargin=1.15em,itemsep=0.12em,topsep=0.2em,parsep=0pt}
\setlist[enumerate]{leftmargin=1.5em,itemsep=0.18em,topsep=0.2em,parsep=0pt}
\setlist[description]{leftmargin=0pt,itemsep=0.16em,topsep=0.2em,parsep=0pt,
  font=\normalfont\bfseries}

% ---------------------------------------------------------------------------
% Document identity. Each leaf declares these before \begin{document} so the
% running head, the colophon, and the PDF metadata all agree.
% ---------------------------------------------------------------------------
\newcommand{\TelosProject}{Telos}
\newcommand{\TelosDocTitle}{}
\newcommand{\TelosDocSubtitle}{}
\newcommand{\TelosRevision}{}

\newcommand{\telosmeta}[4]{%
  \renewcommand{\TelosProject}{#1}%
  \renewcommand{\TelosDocTitle}{#2}%
  \renewcommand{\TelosDocSubtitle}{#3}%
  \renewcommand{\TelosRevision}{#4}%
  \hypersetup{pdftitle={#2},pdfsubject={#3; version #4},
    pdfkeywords={Telos, version #4},pdfauthor={Telos}}%
}

\pagestyle{fancy}
\fancyhf{}
\renewcommand{\headrulewidth}{0pt}
\renewcommand{\footrulewidth}{0.4pt}
\fancyfoot[L]{\footnotesize\textsc{\TelosProject}}
\fancyfoot[C]{\footnotesize\TelosDocTitle}
\fancyfoot[R]{\footnotesize\thepage}

% The first page of a one-pager carries the masthead instead of a running foot,
% so its footer is the compact provenance line.
\fancypagestyle{telosfirst}{%
  \fancyhf{}%
  \renewcommand{\headrulewidth}{0pt}%
  \renewcommand{\footrulewidth}{0.4pt}%
  \fancyfoot[L]{\footnotesize\textsc{\TelosProject}}%
  \fancyfoot[C]{\footnotesize\TelosRevision}%
  \fancyfoot[R]{\footnotesize\thepage}%
}

% ---------------------------------------------------------------------------
% Masthead. A heavy rule, the title, a light subtitle, a thin rule. The
% optional argument is a right-aligned tag (a species' scientific name, a
% lesson number) that sits on the title's baseline.
% ---------------------------------------------------------------------------
\newcommand{\telosmasthead}[1][]{%
  \thispagestyle{telosfirst}%
  \begingroup
  \setlength{\parskip}{0pt}%
  \noindent\rule{\linewidth}{1.6pt}\par
  \vspace{0.30em}
  \noindent
  \begin{minipage}[b]{0.66\linewidth}
    \raggedright{\LARGE\bfseries\TelosDocTitle}
  \end{minipage}%
  \hfill
  \begin{minipage}[b]{0.32\linewidth}
    \raggedleft{\small\itshape #1}
  \end{minipage}\par
  \vspace{0.18em}
  \noindent{\small\TelosDocSubtitle}\par
  \vspace{0.28em}
  \noindent\rule{\linewidth}{0.4pt}\par
  \endgroup
  \vspace{0.35em}
}

% ---------------------------------------------------------------------------
% Headings. Compact, small-caps, rule-anchored — a card has no room for the
% article class's default vertical air.
% ---------------------------------------------------------------------------
\titleformat{\section}{\normalfont\large\bfseries}{}{0pt}{}[\vspace{-0.55em}\rule{\linewidth}{0.4pt}]
\titlespacing*{\section}{0pt}{0.85em}{0.35em}
\titleformat{\subsection}{\normalfont\normalsize\bfseries}{}{0pt}{}
\titlespacing*{\subsection}{0pt}{0.6em}{0.2em}
\titleformat{\subsubsection}{\normalfont\small\bfseries\scshape}{}{0pt}{}
\titlespacing*{\subsubsection}{0pt}{0.45em}{0.12em}

% A short run-in heading for card interiors that must not consume a full line.
\newcommand{\telosrubric}[1]{\textbf{#1}\enspace}

% Cross-reference helper. Every Telos project is a set of loose printed sheets,
% so a reference names the sheet rather than a page number.
\newcommand{\sheetref}[1]{\textit{see the \textbf{#1} sheet}}

% Which decision, ADR or sheet governs the paragraph you are reading. A reader
% who disagrees with an instruction should be able to find the reasoning behind
% it rather than having to take it on trust.
\newcommand{\governedby}[1]{{\footnotesize\textsc{governed by #1}}}

% ---------------------------------------------------------------------------
% Quantities. siunitx is not in this TeX Live split, and the guides need only a
% fixed thin space and upright units, so define that directly rather than take
% a package dependency the documented install would not satisfy.
%   \qty{9}{kV}  ->  9 kV        \unitof{kV}  ->  kV
% ---------------------------------------------------------------------------
\newcommand{\unitof}[1]{\ensuremath{\mathrm{#1}}}
\newcommand{\qty}[2]{#1\,\unitof{#2}}
\newcommand{\numrange}[2]{#1\,--\,#2}
\newcommand{\qtyrange}[3]{#1\,--\,#2\,\unitof{#3}}
% Inches and feet appear constantly in the fishing guide; keep one spelling.
\newcommand{\inch}[1]{#1\,in}
\newcommand{\feet}[1]{#1\,ft}

% ---------------------------------------------------------------------------
% Cards and boxes.
% ---------------------------------------------------------------------------

% Titled card with a hairline frame. The default card for facts a reader looks
% up rather than reads through.
\newtcolorbox{teloscard}[1]{
  enhanced, breakable,
  colback=telospaper, colframe=telosink, colbacktitle=telospaper,
  coltitle=telosink,
  boxrule=0.5pt, sharp corners,
  left=6pt, right=6pt, top=4pt, bottom=4pt,
  toptitle=3pt, bottomtitle=3pt, titlerule=0.4pt,
  before skip=6pt, after skip=6pt,
  title={#1}, fonttitle=\small\bfseries
}

% Untitled left-ruled block: hierarchy without a redundant label.
\newtcolorbox{telosframe}{
  enhanced, breakable,
  colback=telospaper, frame hidden, boxrule=0pt, sharp corners,
  borderline west={1.3pt}{0pt}{telosink},
  left=8pt, right=2pt, top=3pt, bottom=3pt,
  before skip=6pt, after skip=6pt
}

% Filled emphasis panel — the "read this first" block. The 10% fill stays
% legible under a cheap laser toner curve and still reads as distinct.
\newtcolorbox{telospanel}[1]{
  enhanced, breakable,
  colback=telosfill, colframe=telosink, colbacktitle=telosfill,
  coltitle=telosink,
  boxrule=0.5pt, sharp corners,
  left=6pt, right=6pt, top=4pt, bottom=4pt,
  toptitle=3pt, bottomtitle=3pt, titlerule=0.4pt,
  before skip=6pt, after skip=6pt,
  title={#1}, fonttitle=\small\bfseries
}

% Hazard block. Deliberately the heaviest object on any page: a 1.4pt frame
% plus a solid title bar. Nothing else in Telos is allowed to look like this,
% so a reader skimming a lesson cannot mistake it for ordinary emphasis.
\newtcolorbox{teloshazard}[1]{
  enhanced, breakable,
  colback=telospaper, colframe=telosink,
  colbacktitle=telosink, coltitle=telospaper,
  boxrule=1.4pt, sharp corners,
  left=6pt, right=6pt, top=5pt, bottom=5pt,
  toptitle=3pt, bottomtitle=3pt,
  before skip=8pt, after skip=8pt,
  title={\MakeUppercase{#1}}, fonttitle=\small\bfseries
}

% A stop-and-check gate. Used before anything destructive, and before anything
% that would be expensive to discover later. Shared by every project that has
% a procedure someone could get hurt following carelessly.
\newtcolorbox{gate}[1]{
  enhanced, breakable,
  colback=telospaper, colframe=telosink,
  boxrule=1.0pt, sharp corners,
  left=6pt, right=6pt, top=4pt, bottom=4pt,
  toptitle=3pt, bottomtitle=3pt, titlerule=0.4pt,
  before skip=7pt, after skip=7pt,
  colbacktitle=telospaper, coltitle=telosink,
  title={\MakeUppercase{#1}}, fonttitle=\small\bfseries
}

% ---------------------------------------------------------------------------
% Rating dots. A five-step scale that prints identically in black-only: filled
% versus open circles, never a shaded bar.
% ---------------------------------------------------------------------------
\newcommand{\telosdot}{\tikz[baseline=-0.42ex]\fill[telosink] (0,0) circle (0.28ex);}
\newcommand{\telosopendot}{\tikz[baseline=-0.42ex]\draw[telosink,line width=0.28pt] (0,0) circle (0.28ex);}
\makeatletter
\newcounter{telos@rate}
\newcommand{\telosrating}[1]{%
  \begingroup
  \setcounter{telos@rate}{0}%
  \@whilenum\value{telos@rate}<#1\do{\telosdot\kern0.22em\stepcounter{telos@rate}}%
  \@whilenum\value{telos@rate}<5\do{\telosopendot\kern0.22em\stepcounter{telos@rate}}%
  \endgroup
}
\makeatother

% ---------------------------------------------------------------------------
% Tables.
%
% tabularx reads its body by scanning for a literal \end{tabularx}, so it can
% never be wrapped in a \newenvironment. Documents therefore open tabularx and
% longtable directly and take their house style from these column types and the
% booktabs rules. The one exception is the fact strip below, which is built on
% plain tabular precisely so it *can* be an environment.
% ---------------------------------------------------------------------------
\newcolumntype{L}{>{\raggedright\arraybackslash}X}
\newcolumntype{R}{>{\raggedleft\arraybackslash}X}
\newcolumntype{C}{>{\centering\arraybackslash}X}
\newcolumntype{B}{>{\bfseries\raggedright\arraybackslash}X}
\newcolumntype{N}{>{\raggedright\arraybackslash}p{0.16\linewidth}}

% Key/value fact strip. Two flush columns of short lookups; the label column is
% a fixed fraction so a stack of cards aligns across the whole guide.
\newenvironment{telosfacts}[1][0.24]
  {\begingroup\small
   \setlength{\parskip}{0pt}%
   \renewcommand{\arraystretch}{1.25}%
   \begin{tabular}{@{}>{\bfseries\raggedright\arraybackslash}p{\dimexpr#1\linewidth\relax}@{\hspace{0.7em}}>{\raggedright\arraybackslash}p{\dimexpr\linewidth-#1\linewidth-0.7em\relax}@{}}}
  {\end{tabular}\par\endgroup}
\newcommand{\telosfact}[2]{#1 & #2\\}

% ---------------------------------------------------------------------------
% Seasonal / diurnal activity bar.
%
% \telosseasonbar{4,3,2,1,1,2,3,3,4,5,4,3} draws a twelve-cell month strip
% where each value 0-5 is a fill height. Height, not gray value, carries the
% signal, so it is exact under any toner curve; the cell is outlined either way
% so an empty month is still visibly a month.
% ---------------------------------------------------------------------------
\newcommand{\telosseasonlabels}{J,F,M,A,M,J,J,A,S,O,N,D}
\newcommand{\telosseasonbar}[1]{%
  \begingroup
  \begin{tikzpicture}[x=1.25em,y=2.1ex,baseline=0pt]
    \foreach \v [count=\i from 0] in {#1} {
      \draw[telosrule,line width=0.3pt] (\i,0) rectangle (\i+0.86,5);
      \ifnum\v>0
        \fill[telosink] (\i,0) rectangle (\i+0.86,\v);
      \fi
    }
    \foreach \m [count=\i from 0] in \telosseasonlabels {
      \node[font=\tiny,anchor=north,inner sep=1.2pt] at (\i+0.43,-0.1) {\m};
    }
  \end{tikzpicture}%
  \endgroup
}

% Four-cell day-part strip: dawn, midday, dusk, night.
\newcommand{\telosdaylabels}{Dawn,Mid,Dusk,Night}
\newcommand{\telosdaybar}[1]{%
  \begingroup
  \begin{tikzpicture}[x=2.9em,y=2.1ex,baseline=0pt]
    \foreach \v [count=\i from 0] in {#1} {
      \draw[telosrule,line width=0.3pt] (\i,0) rectangle (\i+0.92,5);
      \ifnum\v>0
        \fill[telosink] (\i,0) rectangle (\i+0.92,\v);
      \fi
    }
    \foreach \m [count=\i from 0] in \telosdaylabels {
      \node[font=\tiny,anchor=north,inner sep=1.2pt] at (\i+0.46,-0.1) {\m};
    }
  \end{tikzpicture}%
  \endgroup
}

% A bar needs vertical room for its labels when it sits in a fact strip.
\newcommand{\telosbarrow}[1]{\rule{0pt}{4.2ex}#1\rule[-1.6ex]{0pt}{1.6ex}}

% ---------------------------------------------------------------------------
% Plates. A species or apparatus illustration with a caption rule beneath it.
% \telosplate{width fraction}{file}{caption}
% ---------------------------------------------------------------------------
\newcommand{\telosplate}[3]{%
  \begingroup
  \centering
  \includegraphics[width=#1\linewidth]{#2}\par
  \vspace{0.15em}
  \begin{minipage}{#1\linewidth}
    \rule{\linewidth}{0.4pt}\par
    \vspace{0.1em}
    {\footnotesize #3\par}
  \end{minipage}\par
  \endgroup
  \vspace{0.3em}
}

% TikZ drawing conventions shared by every rig, knot, and circuit diagram.
\tikzset{
  telos line/.style={draw=telosink,line width=0.7pt},
  telos hair/.style={draw=telosink,line width=0.35pt},
  telos heavy/.style={draw=telosink,line width=1.2pt},
  telos node/.style={draw=telosink,fill=telospaper,align=center,inner sep=4pt,font=\small},
  telos emphasis/.style={draw=telosink,line width=1.2pt,fill=telospaper,align=center,
    inner sep=4pt,font=\small\bfseries},
  telos label/.style={font=\scriptsize,inner sep=1.5pt},
  telos arrow/.style={-{Stealth[length=1.8mm]},draw=telosink,line width=0.6pt},
  telos dim/.style={|<->|,draw=telosink,line width=0.35pt},
  telos water/.style={draw=telosrule,line width=0.4pt,decorate,
    decoration={snake,amplitude=0.4mm,segment length=3mm}},
}

% ---------------------------------------------------------------------------
% Lesson furniture
% ---------------------------------------------------------------------------

% \lessonhead{number}{time}{who does what}
\newcommand{\lessonhead}[3]{%
  \par\noindent
  \begin{minipage}{\linewidth}
    \rule{\linewidth}{0.4pt}\par
    \vspace{0.3em}
    {\small\textbf{Lesson #1}\quad$\cdot$\quad #2\quad$\cdot$\quad #3\par}
    \vspace{0.25em}
    \rule{\linewidth}{0.4pt}
  \end{minipage}\par
  \vspace{0.5em}}

% The materials list. Two columns of short lines; a lesson that needs a long
% shopping list is a lesson that has not been thought through.
\newenvironment{materials}
  {\begin{telospanel}{What you need}\begin{multicols}{2}
   \begin{itemize}[leftmargin=1em,itemsep=0.05em,topsep=0pt]}
  {\end{itemize}\end{multicols}\end{telospanel}}

% The four rubrics every lesson carries, in the same order every time:
% what you are about to see, what to do, what it means, and what to ask.
\newcommand{\lessonsection}[1]{\section*{#1}\vspace{-0.35em}\rule{\linewidth}{0.4pt}\vspace{0.2em}}

% A prediction box. The point of the curriculum is that they commit to an
% answer before the demonstration, not after.
\newtcolorbox{predict}{
  enhanced, breakable,
  colback=telospaper, colframe=telosink,
  boxrule=0.5pt, sharp corners,
  left=6pt, right=6pt, top=4pt, bottom=4pt,
  before skip=6pt, after skip=6pt,
  borderline west={2.2pt}{0pt}{telosink},
  title={\small\bfseries Predict first}, fonttitle=\small\bfseries,
  colbacktitle=telospaper, coltitle=telosink,
  toptitle=3pt, bottomtitle=3pt, titlerule=0.4pt,
}

% ---------------------------------------------------------------------------
% Colophon. Rendered once, at the end of the last page, in the recovered footer
% area so it never creates a page of its own.
% ---------------------------------------------------------------------------
\newif\ifTelosColophonRendered
\newcommand{\TelosColophonExtra}{}
\newcommand{\TelosColophon}{%
  \ifTelosColophonRendered\else
  \global\TelosColophonRenderedtrue
  \par
  \enlargethispage{0.5in}%
  \vspace{0.3em}%
  \begingroup
  \setlength{\parskip}{0pt}%
  \begin{minipage}{\linewidth}
  \rule{\linewidth}{0.35pt}\par
  \vspace{0.3em}%
  \fontsize{7}{8.2}\selectfont
  \textbf{Telos.} \TelosProject{} — \TelosDocTitle. Revised \TelosRevision.
  \TelosColophonExtra{}
  Project-created text and diagrams are licensed \href{https://creativecommons.org/licenses/by/4.0/}{CC BY 4.0}.
  Incorporated public-domain artwork remains public domain and is credited on
  the plate it appears with. See \texttt{LICENSE} and \texttt{CREDITS.md} in the
  source. Nothing here is professional advice; verify current regulations and
  follow the stated safety procedure.
  \end{minipage}
  \endgroup
  \fi
}
\AtEndDocument{\TelosColophon}

% Shared macros for the homelab project.
%
% This project's documents are read in two very different situations: at a desk
% while designing, and standing in front of a broken machine at eleven at night
% with no network. The second case is the one that sets the style. Commands are
% set apart so they can be found by eye, every dangerous step has a visible
% verification after it, and nothing depends on being able to reach a web page.

% ---------------------------------------------------------------------------
% Decision-record entries, used by the generated digest.
% ---------------------------------------------------------------------------
\newcommand{\adrentry}[4]{%
  \par\vspace{0.7em}
  \needspace{4\baselineskip}%
  \noindent
  \begin{minipage}{\linewidth}
    \rule{\linewidth}{0.6pt}\par
    \vspace{0.25em}
    {\small\textbf{ADR #1}\quad #2\par}
    \vspace{0.1em}
    {\fontsize{7}{8.2}\selectfont\textsc{#3}\quad$\cdot$\quad #4\par}
    \vspace{0.2em}
    \rule{\linewidth}{0.3pt}
  \end{minipage}\par
  \vspace{0.3em}}

% ---------------------------------------------------------------------------
% Procedures.
%
% A rebuild manual is a sequence of commands and the observable evidence that
% each one worked. These environments keep that pairing structural rather than
% a matter of the author remembering.
% ---------------------------------------------------------------------------

% A block of shell to type. Deliberately not a listing package: plain verbatim
% in a ruled box prints identically everywhere and has no font dependency.
\newenvironment{shellblock}
  {\par\vspace{0.25em}\noindent
   \begin{minipage}{\linewidth}
   \rule{\linewidth}{0.3pt}\par\vspace{0.15em}
   \begingroup\footnotesize\ttfamily\obeylines\obeyspaces\parindent=0pt\parskip=0pt}
  {\endgroup\par\vspace{0.15em}\rule{\linewidth}{0.3pt}\end{minipage}\par\vspace{0.25em}}

% What you should see if the previous step worked. A step without one of these
% is a step you cannot verify, and in a rebuild manual that is a defect.
\newcommand{\expectthis}[1]{%
  \par\vspace{0.1em}
  \noindent\begin{minipage}{\linewidth}
  \hspace*{0.6em}\begin{minipage}{\dimexpr\linewidth-0.6em\relax}
  {\footnotesize\textbf{Expect:} #1\par}
  \end{minipage}\end{minipage}\par\vspace{0.3em}}

% \begin{gate} now lives in common/preamble.tex, shared with the other
% projects that document procedures.

% \governedby now lives in common/preamble.tex.

% An instance value that comes from the private overlay. In the published
% documents these render as visible placeholders, which is deliberate: a
% published manual must be complete and obviously generic, never look like it
% is describing somebody's actual network.
\newcommand{\instancevalue}[2]{\texttt{\textlangle #1\textrangle}%
  \footnote{Instance value. Supplied from \texttt{homelab/instance/} at build
  time; #2.}}
\newcommand{\placeholder}[1]{\texttt{\textlangle #1\textrangle}}

% ---------------------------------------------------------------------------
% Role and boundary diagrams.
% ---------------------------------------------------------------------------
\tikzset{
  role/.style={draw=telosink,line width=0.9pt,fill=telospaper,align=center,
    inner sep=5pt,font=\small,minimum width=2.6cm},
  roleoff/.style={draw=telosmute,line width=0.5pt,dash pattern=on 2.5pt off 1.8pt,
    fill=telospaper,align=center,inner sep=5pt,font=\small,minimum width=2.6cm},
  boundary/.style={draw=telosmute,line width=0.5pt,dash pattern=on 3pt off 2pt,
    rounded corners=3pt},
  flow/.style={-{Stealth[length=2mm]},draw=telosink,line width=0.7pt},
  hlabel/.style={font=\fontsize{6.4}{7.4}\selectfont},
}

% GENERATED by scripts/adr-digest from homelab/decisions/*.md.
% Edit the Markdown and re-run; do not edit this file.
\telosmeta{Homelab}{Decision Record}{Every accepted, superseded and deferred decision, in order}{2026-07-25}
\begin{document}
\telosmasthead[62 decisions]

This is the complete architecture decision record for the homelab provisioning system, generated from the Markdown sources so the printed copy cannot drift from the repository. Each entry carries its status and its Decision section verbatim. Context and consequences remain in \texttt{homelab/decisions/}.

\begin{telosfacts}[0.30]
\telosfact{Accepted}{44}
\telosfact{Deferred for post-proof revalidation}{14}
\telosfact{Superseded}{4}
\end{telosfacts}

\adrentry{0001}{Local repository location}{Superseded by ADR 0046}{2026-07-24}
Use \texttt{/home/ksh/git/codex/homelab-infrastructure} as the normal local working
copy and initialize it as a Git repository with \texttt{main} as its initial branch.

\adrentry{0002}{Initial computer profiles}{Accepted}{2026-07-24}
Begin with two computer profiles:
\begin{itemize}
\item \textbf{Controller}, identified as \texttt{controller}
\item \textbf{Workstation}, identified as \texttt{workstation}
\end{itemize}
The Controller is the homelab server, orchestrator, and general infrastructure
utility. Each Workstation deployment pivots on the hostname of its Controller.

\adrentry{0003}{Design-phase commit granularity}{Accepted}{2026-07-24}
During the initial architecture and design phase, accumulate and integrate
related decisions into larger coherent sections before committing them.
After the design foundation exists, use smaller, focused commits for
implementation work and later revisions.

\adrentry{0004}{Interactive install authorization}{Accepted}{2026-07-24}
Do not require a pre-staged deployment record. The network-boot environment
will collect and validate several operator inputs at the target computer,
display a complete preflight summary, and perform no destructive operation
until the operator gives a final explicit confirmation.
Physical-console and boot-path access, together with this interactive
confirmation workflow, are the authorization boundary.

\adrentry{0005}{Use home.arpa for internal DNS}{Accepted}{2026-07-24}
Use \texttt{home.arpa} as the internal DNS suffix. Controller references use a fully
qualified hostname beneath it, such as \texttt{polycarp.home.arpa}.

\adrentry{0006}{Keep baseline and bootstrap DNS on UniFi}{Superseded by ADR 0008}{2026-07-24}
Keep UniFi as the normal client-facing DNS endpoint and the owner of the
Controller's bootstrap record. For a Controller whose instance hostname is
\texttt{polycarp}, that record is \texttt{polycarp.home.arpa}.
Defer Controller-hosted DNS. A later design may place an authoritative child
zone on the Controller and configure a UniFi Forward Domain rule for that child
zone, but it must preserve the UniFi-hosted Controller bootstrap record.

\adrentry{0007}{Design generic profiles and build a Controller first}{Accepted}{2026-07-24}
Design reusable Controller and Workstation profiles without using an inventory
of the current Polycarp installation as an architecture input.
The first end-to-end acceptance test will network boot a target and use the
generated provisioning system to build a Controller from bare metal. The
acceptance sequence will then validate Controller-owned DHCP and DNS on an
isolated network whose only permanent network infrastructure is a switch.

\adrentry{0008}{Support Controller-owned DHCP and DNS}{Accepted}{2026-07-24}
The Controller must optionally function as the primary DHCP and DNS server for
its managed network.
When external infrastructure owns DHCP or DNS, the corresponding Controller
services must remain disabled. Two DHCP authorities must never be active on the
same layer-2 network during installation, testing, or transition.
The first acceptance sequence must validate the installed Controller's DHCP and
DNS behavior on an isolated network containing no permanent router or DNS/DHCP
appliance. Provisioning occurs first on the existing full network, and ADR 0010
defines the powered-off transition to the isolated network.

\adrentry{0009}{Activate Controller network services automatically}{Accepted}{2026-07-24}
When the Controller is provisioned for Controller-owned DHCP and DNS, activate
those services automatically on its first boot after provisioning. Do not
require a separate manual activation command.
The activation path must fail closed: if its required transition conditions
cannot be verified, Controller DHCP must remain stopped.

\adrentry{0010}{Power off before the network-service first boot}{Accepted}{2026-07-24}
When Controller-owned DHCP and DNS are selected, a successful installation
must end with a clean poweroff rather than rebooting into the installed
operating system.
The operator relocates the powered-off Controller to the isolated switch and
powers it on. That power-on is the first boot of the installed system, and ADR
0009 automatically activates the configured network services after validation.

\adrentry{0011}{Exclude routing and NAT from the initial Controller}{Accepted}{2026-07-24}
Exclude routing, NAT, and perimeter-firewall duties from the initial Controller
profile.
On the isolated acceptance network, the Controller provides local DHCP and DNS
without advertising or acting as a default gateway. The initial isolated test
therefore has no Internet access.
ADR 0045 subsequently excludes a gateway from the network input schema and
derives the absence of the DHCP default-router option.
Reserve a clearly marked future-extension section in the Controller design for
possible routing and NAT support.

\adrentry{0012}{Bundle Controller DHCP and DNS}{Accepted}{2026-07-24}
Expose one combined Controller network-services option in the initial profile:
\begin{itemize}
\item When enabled, Controller DHCP and DNS are configured and activated together.
\item When disabled, both remain stopped and external infrastructure provides those
functions.
\end{itemize}
Reserve independent DHCP-only and DNS-only operation as a future extension.
ADR 0044 subsequently selects one host-native dnsmasq instance as the initial
implementation of this bundle.

\adrentry{0013}{Start Controller network services with IPv4}{Accepted}{2026-07-24}
Limit the initial bundled Controller network-services option to managed IPv4:
\begin{itemize}
\item provide DHCPv4;
\item serve DNS on the Controller's IPv4 service address;
\item do not provide DHCPv6, router advertisements, or managed IPv6 addressing.
\end{itemize}
Keep the operating system's IPv6 stack and ordinary link-local behavior
enabled. Reserve managed dual-stack operation as a future extension.

\adrentry{0014}{Use a host-first hybrid Controller}{Accepted}{2026-07-24}
Use a host-first hybrid architecture:
\begin{itemize}
\item The Controller host operating system directly owns boot, storage, physical networking,
first-boot safety, bundled DHCP/DNS, PXE control, and local recovery.
\item Use containers when they materially simplify packaging or independent
service lifecycle.
\item Use VMs only when a service requires a separate kernel, different operating
system, or stronger trust boundary.
\item Do not require virtualization for every service.
\end{itemize}

\adrentry{0015}{Use Arch Linux for the Controller host}{Accepted}{2026-07-24}
Use Arch Linux as the host operating system for the initial Controller profile.
This decision selected the operating-system family only. Subsequent ADRs select
the kernels and Btrfs-inside-LUKS2 root filesystem. The installation package
state and complete rollback mechanism remain open.
FreeBSD is not part of the initial Controller baseline. It remains available
for a future profile or for a workload that specifically benefits from it.

\adrentry{0016}{Use guarded direct Arch updates}{Accepted}{2026-07-24}
Use pacman directly for guarded Controller updates:
\begin{itemize}
\item do not apply unattended operating-system package upgrades;
\item review Arch news and the complete pending update set before maintenance;
\item create a recoverable system checkpoint before changing packages;
\item perform a supported full-system pacman upgrade, not a partial upgrade;
\item reboot when the transaction affects the kernel, initramfs, system manager,
networking stack, or foundational services;
\item automatically validate Controller health, including DHCP, DNS, and PXE,
after the update; and
\item restore the checkpoint or rebuild from recorded inputs if validation fails.
\end{itemize}
Do not create a private Arch repository, a package-promotion service, or a
wrapper package manager for the initial Controller.
This ADR defines the workflow boundary, not its implementation. Subsequent ADRs
select the kernels and Btrfs-inside-LUKS2 root filesystem, and ADR 0028 selects
native Btrfs checkpoint operations. ADR 0032 requires the off-host Secure Boot
signing path to pass preflight before a hardened full-system upgrade mutates
packages. ADR 0043 defers that signing requirement during the functional proof:
the proof workflow must instead regenerate, checksum, install, and validate the
matching development UKIs as part of the guarded transaction. Exact health
checks, maintenance cadence, retention, and package-state recording format
remain undecided.

\adrentry{0017}{Use linux-lts as the default Controller kernel}{Accepted}{2026-07-24}
Use Arch's official \texttt{linux-lts} package as the default kernel for the initial
Controller profile.
This decision did not determine whether the profile also installs a secondary
kernel; ADR 0018 subsequently selected the standard \texttt{linux} package for that
role. This decision does not alter the guarded full-system update policy in ADR
0016.

\adrentry{0018}{Install the standard Linux kernel as a secondary kernel}{Accepted}{2026-07-24}
Install Arch's official standard \texttt{linux} package as the secondary Controller
kernel. Keep \texttt{linux-lts} as the default boot selection.
Both kernels are managed by pacman and updated in the same guarded full-system
transaction defined by ADR 0016.
ADR 0022 subsequently selects systemd-boot. The eventual configuration must
expose an operator-selectable standard-kernel entry without making it the
automatic default; exact menu behavior remains open.

\adrentry{0019}{Require UEFI boot for the Controller}{Accepted}{2026-07-24}
Require the initial Controller profile to provision and boot in UEFI mode.
Legacy BIOS and Compatibility Support Module boot are unsupported.
The installer preflight must verify that it is running in UEFI mode before
offering destructive authorization. The eventual installed-system acceptance
test must verify UEFI boot with \texttt{linux-lts} as the default and standard \texttt{linux}
available as the secondary kernel.
This decision applies to the Controller profile. It does not independently
decide the firmware requirements for every future Workstation target.
ADR 0020 subsequently requires Secure Boot and encryption at rest. ADR 0030
retains factory PK and KEK ownership while adding the lab signer to \texttt{db}.
Disk partitioning and the precise UEFI network-boot chain remain separate
decisions. ADR 0022 subsequently selects systemd-boot for Arch, and ADR 0023
selects signed Unified Kernel Images.

\adrentry{0020}{Require encryption at rest and authenticated boot}{Accepted}{2026-07-24}
For systems installed and managed by the Controller and Workstation profiles:
\begin{itemize}
\item encrypt every persistent data-bearing operating-system and service volume;
\item encrypt disk-backed swap and persistent crash or hibernation state;
\item encrypt backups of managed-system data, secrets, and recovery state;
\item use LUKS2 for Linux data-bearing volumes;
\item use BitLocker for Windows data-bearing volumes;
\item require Secure Boot for provisioning, installed-system boot, and recovery
paths;
\item limit unencrypted storage to mandatory firmware-readable boot and recovery
partitions;
\item keep those unencrypted partitions minimal and free of credentials,
decryption keys, tokens, private data, and other reusable secrets; and
\item authenticate executable boot artifacts through the accepted Secure Boot
chain.
\end{itemize}
This is an at-rest encryption boundary. It does not yet define transport
encryption.
The policy applies to new or rebuilt systems managed by these profiles. It does
not assert that existing lab devices have already been migrated or authorize
inventorying them.
ADR 0043 creates one narrow sequencing exception. Its labeled, disposable
\texttt{development-proof} installation must still use encryption at rest, but may
defer the custom authenticated-boot chain and run with Secure Boot disabled
when an already trusted upstream chain is unavailable. Such an installation
does not satisfy this ADR for final profile acceptance and must be reprovisioned
after the signing design is revisited.

\adrentry{0021}{Install a permanent Windows maintenance environment}{Superseded by ADR 0047}{2026-07-24}
Include a permanent, licensed, bare-metal Windows maintenance installation in
the initial Controller profile.
\begin{itemize}
\item Use Windows only for physical firmware and hardware maintenance.
\item Keep Arch Linux as the default boot and the sole normal Controller runtime.
\item Do not host Controller services or authoritative configuration in Windows.
\item Protect the Windows operating-system volume with BitLocker.
\item Keep Windows Boot Manager independently registered as a native UEFI boot
option.
\item Store the BitLocker recovery material off-host under the recovery policy
required by ADR 0020.
\item Treat WinPE and vendor-supplied boot media as supplemental recovery or
one-off tools, not as the supported primary Windows maintenance path.
\end{itemize}
The Windows edition, license source, storage allocation, account model, update
cadence, direct-boot user experience, and exact relationship between Windows
Boot Manager and the Linux boot manager remained undecided here. ADR 0022
subsequently selects systemd-boot for Arch while preserving an independent
native Windows UEFI entry; the direct-boot user experience remains open.
ADR 0043 permits the Windows installation and direct UEFI entry to be exercised
during the functional proof. That test is not final Windows acceptance;
BitLocker, Secure Boot, TPM, and recovery behavior must be completed and
retested after boot hardening.

\adrentry{0022}{Use systemd-boot for Arch}{Accepted}{2026-07-24}
Use systemd-boot as the installed Arch boot manager for the initial Controller
profile.
\begin{itemize}
\item Make the systemd-boot UEFI entry the normal firmware default.
\item Make \texttt{linux-lts} the automatic Arch kernel selection.
\item Expose standard \texttt{linux} as an operator-selectable secondary Arch entry.
\item Keep Windows Boot Manager independently registered as a native UEFI entry.
\item Do not make Windows depend on systemd-boot remaining functional.
\item Sign and trust systemd-boot as part of the Secure Boot chain required by ADR
0020.
\end{itemize}
The operator experience for requesting a Windows boot remains undecided.
systemd-boot may discover Windows, but the final design must retain a direct
firmware path that does not require chainloading through the Linux boot manager.
ADR 0023 subsequently selects signed Unified Kernel Images as the kernel
artifact format. ADR 0030 retains factory PK and KEK ownership while adding the
lab signer to \texttt{db}, and ADR 0031 makes that signer lab-wide. Private-key
custody is subsequently resolved by ADR 0032. Signing automation, boot
counting, checkpoint integration, and boot-menu timeout remain undecided. ADR
0024 subsequently adds a self-contained recovery UKI, and ADR 0025 places all
boot artifacts on one shared ESP without XBOOTLDR.
Under ADR 0043, Milestone A may exercise the same systemd-boot and UKI topology
without the custom trust and signing workflow. Signing remains mandatory for
final acceptance after the deferred design is revalidated.

\adrentry{0023}{Use signed Unified Kernel Images}{Accepted}{2026-07-24}
Build and install a signed UKI for each accepted Arch kernel:
\begin{itemize}
\item one UKI for \texttt{linux-lts}, selected by default; and
\item one UKI for standard \texttt{linux}, exposed as the secondary Arch entry.
\end{itemize}
Each UKI contains its matching kernel, initramfs, and kernel command line.
systemd-boot discovers and launches the UKIs as Boot Loader Specification Type
\#2 entries.
Do not make an unsigned or partially authenticated Linux boot entry part of the
normal installed-system path. The Secure Boot trust chain must verify each UKI
before execution.
UKIs reside on firmware-readable, unencrypted boot storage and are therefore
public artifacts. They must not contain raw volume keys, credentials, reusable
tokens, private data, or other secrets.
This decision does not select the UKI generator, initramfs generator, signing
tool, or signing automation. ADR 0024 subsequently selects a third
self-contained recovery UKI, ADR 0025 selects one shared ESP for all boot
artifacts, ADR 0029 defines TPM2 unlock authorization for the two normal UKIs,
ADR 0030 defines additive enrollment of the lab signer without replacing
factory PK or KEK ownership, and ADR 0032 defines removable off-host custody of
the signing private key.
ADR 0043 keeps the UKI topology but allows checksum-verified, unsigned
development UKIs only in the explicitly non-production functional proof. That
exception is removed before final acceptance; it is not an installed profile
option.

\adrentry{0024}{Include a self-contained recovery UKI}{Accepted}{2026-07-24}
Install a third, signed, self-contained recovery UKI on the Controller.
The recovery UKI:
\begin{itemize}
\item boots independently of the installed Arch root filesystem;
\item is operator-selectable and never the automatic normal boot;
\item contains the storage, encryption, filesystem, and boot-artifact tools needed
to diagnose and recover the installed system;
\item requires separately held LUKS2 recovery material under ADR 0029 and contains
no raw unlock key, credential, token, or other reusable secret;
\item starts in a non-destructive state and requires explicit operator action
before mounting writable storage or changing disk state; and
\item is authenticated by the same accepted Secure Boot trust policy as the normal
Arch UKIs.
\end{itemize}
The exact rescue userspace, image generator, included hardware support,
network behavior, update cadence, validation procedure, and whether firmware
also registers the recovery UKI directly remain undecided.
ADR 0043 permits a checksum-verified, unsigned recovery-UKI candidate during
the functional proof so its userspace and recovery behavior can be exercised.
Its authenticated execution, rotation, and interaction with TPM policy remain
final-hardening acceptance tests.

\adrentry{0025}{Use one shared EFI System Partition}{Accepted}{2026-07-24}
Use one FAT32 EFI System Partition shared by Windows and Arch on the initial
Controller.
\begin{itemize}
\item Store Windows Boot Manager in its standard vendor directory.
\item Store systemd-boot in its standard vendor directory.
\item Store the two normal Arch UKIs and the recovery UKI in the Boot Loader
Specification location on the ESP.
\item Do not create an XBOOTLDR partition.
\item Keep the ESP unencrypted, minimal, secret-free, and covered by Secure Boot
integrity verification under ADR 0020.
\item Back up and validate the complete ESP as one unit.
\end{itemize}
ADR 0043 preserves this physical layout during the functional proof but defers
Secure Boot integrity acceptance. Proof-stage ESP backups and artifacts use
checksums and explicit \texttt{development-proof} labeling until hardening.
The exact ESP size must be selected only after measuring the generated UKIs and
allowing documented update and recovery headroom. Fallback-path ownership,
redundant boot-media copies, backup format, and restore procedure remain
undecided.

\adrentry{0026}{Use Btrfs inside LUKS2 for the Arch root}{Accepted}{2026-07-24}
Use Btrfs for the initial Controller's Arch root filesystem, contained inside a
LUKS2-encrypted volume.
This decision selects the root filesystem and encryption layering only. ADR
0027 subsequently defines the operating-system checkpoint boundary. The
following remain open:
\begin{itemize}
\item compression and mount options;
\item swap implementation;
\item service-data placement and application-consistent backup behavior;
\item multi-device layout or redundancy;
\item storage-capacity allocations.
\end{itemize}
ADR 0028 subsequently selects native Btrfs operations for checkpoints.
ADR 0029 subsequently defines automatic TPM2 root unlock plus an off-host
recovery key.
Additional persistent data-bearing volumes must still satisfy ADR 0020 but do
not automatically inherit this exact filesystem choice.

\adrentry{0027}{Limit root checkpoints to operating-system state}{Accepted}{2026-07-24}
Make the root Btrfs checkpoint an operating-system checkpoint.
The checkpointed root includes:
\begin{itemize}
\item installed operating-system files and libraries;
\item \texttt{/etc} and other declarative host configuration;
\item pacman's database and other package-manager state; and
\item system state that must remain version-aligned with the installed packages.
\end{itemize}
Use separate Btrfs subvolumes, excluded from a root checkpoint, for:
\begin{itemize}
\item ordinary user home directories and the root operator's home;
\item logs;
\item caches and disposable temporary state;
\item the snapshot store itself; and
\item mutable application and Controller service data.
\end{itemize}
Do not exclude all of \texttt{/var/lib} as one unit. Keep pacman and other
version-coupled operating-system state with the root, and assign mutable service
paths to separate subvolumes when each service is defined.
The exact subvolume names, mount options, retention rules, and service paths
remain undecided. ADR 0028 subsequently selects native Btrfs operations as the
checkpoint primitive.

\adrentry{0028}{Manage checkpoints with native Btrfs operations}{Accepted}{2026-07-24}
Use native Btrfs subvolume operations as the checkpoint primitive in the
repository-controlled guarded-update workflow.
\begin{itemize}
\item Create a read-only snapshot of the operating-system root before a guarded
update.
\item Associate that root snapshot with a backup and manifest of the matching
shared-ESP artifacts.
\item Record enough update and package-state metadata to identify the checkpoint
and the transaction it protects.
\item Use the self-contained recovery UKI to restore a root checkpoint and its
matching boot artifacts when the installed root is not safely recoverable.
\item Do not install Snapper, \texttt{snap-pac}, or an automatic timeline-snapshot service
for the initial Controller.
\end{itemize}
This decision selects the checkpoint primitive rather than the complete
workflow. Snapshot naming, metadata format, retention limits, deletion policy,
ESP backup format, exact rollback commands, and post-validation checkpoint
lifecycle remain undecided.

\adrentry{0029}{Automatically unlock the Controller root with TPM2}{Deferred for post-proof revalidation by ADR 0043}{2026-07-24}
ADR 0043 defers this decision during the functional proof. Do not enroll a TPM
unlock token or policy in Milestone A; unlock LUKS2 manually. Revisit this ADR
with ADRs 0030 through 0042 after the functional environment works.
Require TPM 2.0 for the initial Controller and automatically unlock its Arch
LUKS2 root only from either normal signed UKI:
\begin{itemize}
\item the default \texttt{linux-lts} UKI; or
\item the secondary standard \texttt{linux} UKI.
\end{itemize}
Bind the TPM enrollment to a signed policy over the normal UKIs' measured boot
state, using the systemd UKI measurement in PCR 11. Do not rely on an
unconstrained TPM enrollment or only on the current raw PCR values. The exact
additional PCR bindings remain undecided. ADR 0033 subsequently selects a
dedicated TPM-policy signing identity distinct from the Secure Boot signer and
applies ADR 0032's removable off-host custody pattern to its private key. ADR
0034 subsequently shares that policy identity across Controllers without
sharing their TPM-sealed LUKS2 secrets.
Enroll a separate, high-entropy LUKS2 recovery key. Store it off-host and never
in Git, the ESP, a UKI, ordinary logs, or the Controller's only recovery copy.
If TPM policy validation fails, early boot must fall back to recovery-key entry
without weakening or replacing the policy automatically.
Authenticate the recovery UKI with Secure Boot, but do not give it the signed
PCR authorization needed for automatic LUKS2 unlock. Recovery boot always
requires separately held recovery material.
Do not require a TPM PIN for normal Controller boot.
ADR 0030 subsequently retains OEM and Microsoft firmware trust while adding a
lab signer to Secure Boot \texttt{db}. That broader permission to execute an EFI image
does not broaden this TPM policy: only the two approved normal lab UKI measured
states receive automatic root unlock. ADR 0031's lab-wide Secure Boot signing
identity is not thereby selected as the TPM-policy signing identity; ADR 0033
subsequently requires a separate key pair.

\adrentry{0030}{Add lab Secure Boot trust without replacing platform ownership}{Deferred for post-proof revalidation by ADR 0043}{2026-07-24}
Use an additive Secure Boot trust model for the initial Controller profile.
\begin{itemize}
\item Retain the platform's factory or OEM ownership represented by its PK and
KEKs. Do not replace them with lab-owned PK or KEK material.
\item Preserve the compatible OEM and Microsoft trust required for Windows,
firmware servicing, third-party UEFI applications, and option ROMs.
\item Preserve the platform's supported \texttt{dbx} revocation-update path.
\item Add a lab-controlled public signing certificate to the firmware's authorized
signature database (\texttt{db}) through a firmware-supported, physical-presence
enrollment procedure that preserves the existing trust entries.
\item Sign systemd-boot and all three Controller UKIs with the corresponding lab
signing identity.
\item Continue to authorize automatic LUKS2 unlock only for the two normal
lab-approved UKIs under ADR 0029. Firmware permission to execute another
OEM- or Microsoft-trusted EFI image does not grant that image the TPM unlock
capability.
\end{itemize}
Before destructive installation, preflight must establish that the target can
retain the required factory trust and enroll the lab certificate without
replacing it. If that cannot be demonstrated, the initial Controller profile
does not support the target and installation must stop before wipe
authorization.
Record the enrolled certificate fingerprints and the retained trust state in
the machine's non-secret provisioning record. Do not put private signing
material in firmware, the ESP, a UKI, or Git.
ADR 0031 subsequently selects one lab-wide Secure Boot signing identity. The
private-key custody model is subsequently resolved by ADR 0032. Signing
automation, rotation and revocation procedure, exact enrollment tooling, and
current certificate manifest remain separate decisions. The manifest must
account for vendor and Microsoft certificate transitions rather than
permanently assuming the factory's original certificate set is sufficient.

\adrentry{0031}{Use one lab-wide Secure Boot signing identity}{Deferred for post-proof revalidation by ADR 0043}{2026-07-24}
Use one dedicated, lab-wide Secure Boot code-signing identity for managed
Controller and Workstation systems.
\begin{itemize}
\item Enroll the same public certificate in firmware \texttt{db} on every managed machine
that must execute lab-signed EFI artifacts.
\item Use its private counterpart to sign the lab-controlled installed-system EFI
artifacts defined by each profile, including the Controller's systemd-boot
executable and normal and recovery UKIs.
\item Record and verify the shared certificate fingerprint in the profile's
non-secret artifact manifest and each machine's provisioning record.
\item Do not create or manage a separate Secure Boot signing identity for every
machine.
\end{itemize}
Sharing the identity does not mean copying its private key to every machine.
ADR 0020 keeps private signing material off-host, and ADR 0032 subsequently
selects encrypted removable primary and backup media used only in a designated
signing environment. Signing automation, detailed rotation, and operator
authorization remain separate decisions.
This identity is scoped to Secure Boot code signing. This decision does not
authorize its use for TLS, SSH, package repositories, backup encryption, disk
recovery, or other unrelated purposes. ADR 0033 subsequently selects a
different identity for signing TPM PCR policies. The network-provisioning trust
chain remains open.

\adrentry{0032}{Keep the Secure Boot private key on removable off-host media}{Deferred for post-proof revalidation by ADR 0043}{2026-07-24}
Keep the lab-wide Secure Boot private key on a primary encrypted removable
medium that is normally detached and powered off.
Use that medium only in a designated trusted signing environment during planned
provisioning and update work. The existing bootstrap machine may provide the
initial signing environment, but the key must not be mounted in its ordinary
PXE-serving context. The exact isolated or ephemeral signing environment
remains an implementation decision. If that bootstrap machine is later
provisioned as a managed Controller or Workstation, signing responsibility must
move elsewhere first.
Do not copy or persist the private key on:
\begin{itemize}
\item a Controller or Workstation;
\item an EFI System Partition or UKI;
\item a PXE root, installer, or target machine;
\item the ordinary filesystem of the bootstrap host;
\item an always-online signing service; or
\item Git, logs, build output, or provisioning records.
\end{itemize}
Build or prepare artifacts outside the key medium, validate the artifact and
its non-secret manifest in the signing environment, sign it there, verify the
result against the expected public-certificate fingerprint, and return only
the signed artifact and non-secret metadata. Unsigned candidates must remain
outside the active ESP. Package or boot-manager hooks must not install an
unsigned or locally self-signed replacement into the active boot path.
Create one separately encrypted backup copy on a second removable medium. Keep
it offline and physically separate from the primary medium and from the
credential that unlocks either copy. Do not connect the primary and backup
media during the same routine signing session. Periodically verify that the
backup can be decrypted and can produce a signature accepted by the recorded
public certificate; the exact test cadence remains open.
Every guarded Controller full-system upgrade must preflight the signing path
before package mutation, because the transaction may require new systemd-boot,
kernel, initramfs, or UKI signatures. Absence or failure of the signer stops the
update before intentional changes begin. A boot-affecting update is not
successful until all returned artifacts pass signature and manifest validation
and are installed together. Normal boot and ordinary runtime do not require
access to the private key.
ADR 0036 subsequently selects manually unlocked LUKS2 volumes for both
removable media. The signing environment, transfer format, tooling,
temporary-data handling, and detailed rotation and revocation procedure remain
separate decisions. ADR 0040 subsequently selects the operator's existing
off-Controller password manager for unlock-credential custody.
ADR 0033 subsequently applies this custody pattern to a distinct TPM-policy
signing private key. ADR 0035 subsequently places both distinct private keys on
the same physical primary medium and both backup copies on the same physical
backup medium.

\adrentry{0033}{Separate TPM policy signing from Secure Boot signing}{Deferred for post-proof revalidation by ADR 0043}{2026-07-24}
Create a dedicated TPM signed-PCR policy key pair that is cryptographically
distinct from the lab-wide Secure Boot code-signing key pair.
Use the TPM policy private key only to sign approved PCR 11 policies for the two
normal Controller UKIs:
\begin{itemize}
\item the default \texttt{linux-lts} UKI; and
\item the secondary standard \texttt{linux} UKI.
\end{itemize}
Do not give the recovery UKI a policy signature that authorizes automatic root
unlock. Continue to sign it only for Secure Boot execution under ADR 0024.
Enroll or otherwise bind only the TPM policy public key to the Controller's
LUKS2 TPM2 unlock policy. The public key and its fingerprint are non-secret and
belong in the artifact manifest and machine provisioning record. Never enroll,
embed, or copy the TPM policy private key to a managed machine.
Apply ADR 0032's custody pattern to the TPM policy private key: keep encrypted
primary and backup copies on removable off-host media and use the key only in
the designated signing environment during planned work. ADR 0035 subsequently
co-locates both distinct private keys on the same primary medium and both
backup copies on the same backup medium.
For each normal UKI, the returned artifact must contain a valid TPM policy
signature and a valid Secure Boot signature. The final Secure Boot signature
must authenticate the UKI containing its policy metadata. Verify both
authorities before installing the artifact in the active ESP.
ADR 0034 subsequently makes the TPM policy identity shared across Controllers.
ADR 0035 subsequently resolves its physical-media placement. Its algorithm,
lifetime, exact systemd enrollment and signing commands, rotation procedure,
and additional PCR bindings remain separate decisions.

\adrentry{0034}{Share one TPM policy signing identity across Controllers}{Deferred for post-proof revalidation by ADR 0043}{2026-07-24}
Use one dedicated TPM policy signing identity across all machines provisioned
with the Controller profile.
\begin{itemize}
\item Bind the same TPM policy public key to each Controller's LUKS2 signed-PCR
policy.
\item Use the corresponding shared private key to sign approved PCR 11 policies
for each Controller's two normal UKIs, including instance-specific UKIs.
\item Record and verify the common policy-public-key fingerprint in the Controller
profile artifact manifest and each instance's non-secret provisioning
record.
\item Do not generate or maintain a separate TPM policy private key for every
Controller.
\end{itemize}
This does not create a shared LUKS2 recovery key, volume key, TPM-sealed secret,
or TPM enrollment. Those remain unique to each Controller. Moving an encrypted
disk to another Controller does not move its automatic-unlock capability.
The identity remains cryptographically separate from the lab-wide Secure Boot
signer under ADR 0033 and follows ADR 0032's removable off-host custody model.
ADR 0035 subsequently places both private keys on the same primary removable
medium and both backup copies on the same backup medium.
This decision applies only to the accepted Controller automatic-unlock design.
It does not enroll Workstations or future profiles into the policy without a
separate decision.

\adrentry{0035}{Co-locate both signing keys on each removable medium}{Deferred for post-proof revalidation by ADR 0043}{2026-07-24}
Store both distinct private keys on the same encrypted primary removable
medium:
\begin{itemize}
\item the lab-wide Secure Boot code-signing private key; and
\item the Controller-wide TPM-policy signing private key.
\end{itemize}
Store the backup copy of both keys together on the same separately encrypted
backup medium defined by ADR 0032.
Keep the identities logically and cryptographically separate. Give each key a
role-specific name, expected public-key fingerprint, and explicit signing
configuration. Signing tooling must select the intended key by role rather than
expose one ambiguous generic signer.
Rotating either identity must update and verify both the primary and backup
media before the rotation is complete. The other identity does not rotate
merely because the two share storage.
ADR 0036 subsequently selects LUKS2 for both removable media. The filesystem,
subsequently selected by ADR 0037, is ext4. The directory layout, file
encryption and permissions, and detailed backup-synchronization procedure
remain separate decisions. ADR 0039 subsequently selects one shared passphrase
for both media, and ADR 0040 places its authoritative custody in the operator's
existing off-Controller password manager.

\adrentry{0036}{Encrypt signing media with LUKS2}{Deferred for post-proof revalidation by ADR 0043}{2026-07-24}
Use one LUKS2-encrypted data volume on each signing medium:
\begin{itemize}
\item one independently encrypted LUKS2 volume on the primary medium; and
\item one independently encrypted LUKS2 volume on the backup medium.
\end{itemize}
Place all private signing material on the filesystem inside that encrypted
volume. Do not create an unencrypted data filesystem or private-key copy
elsewhere on either medium. The partition table and LUKS2 header are mandatory
unencrypted metadata and must contain no secrets beyond the encrypted keyslot
material intrinsic to LUKS2.
Require an operator-supplied unlock credential in the designated signing
environment. Do not configure TPM auto-unlock, host-stored keyfiles, or another
automatic unlock path for either removable volume.
After signing, unmount the filesystem, close the LUKS mapping, and physically
detach the medium. The signing workflow must fail closed if either operation
does not complete cleanly.
Record the non-secret volume identifiers and the exact cryptographic and key
derivation parameters used to create each volume. Do not expose an unlock
credential in that record.
ADR 0037 subsequently selects ext4 as the inner filesystem. The partition
layout, LUKS2 cipher and PBKDF parameters, LUKS header backup procedure,
file-level key encryption, and rekey procedure remain separate decisions. ADR
0038 subsequently selects the routine mount policy, and ADR 0039 selects one
shared manually entered passphrase. ADR 0040 subsequently places its
authoritative custody in the operator's existing off-Controller password
manager.

\adrentry{0037}{Use ext4 inside the signing-media LUKS2 volumes}{Deferred for post-proof revalidation by ADR 0043}{2026-07-24}
Create one ext4 filesystem inside each signing-media LUKS2 volume:
\begin{itemize}
\item one independently created ext4 filesystem on the primary medium; and
\item one independently created ext4 filesystem on the backup medium.
\end{itemize}
Store private keys, public certificates, fingerprints, and non-secret
role-specific metadata as ordinary files with Unix ownership and permissions.
Do not use FAT, exFAT, Btrfs, ZFS, or another inner filesystem for the initial
signing-media design.
Give the two filesystems distinct, non-secret UUIDs and labels so tooling can
verify that the operator attached the intended primary or backup medium before
opening or synchronizing it. Do not use a device path such as \texttt{/dev/sdX} as
identity.
The exact labels, filesystem creation options, directory layout, ownership,
permissions, mount options, free-space threshold, check cadence, and
backup-synchronization procedure remain separate decisions.
ADR 0038 subsequently selects read-only routine mounts with
\texttt{ro,noload,nodev,nosuid,noexec} and reserves writable mounts for explicit
signing-media maintenance.

\adrentry{0038}{Mount signing media read-only during routine signing}{Deferred for post-proof revalidation by ADR 0043}{2026-07-24}
During routine signing:
\begin{itemize}
\item open the LUKS2 mapping read-only;
\item verify the expected LUKS2 and ext4 identities;
\item require the ext4 filesystem to be clean before mounting it;
\item mount it with \texttt{ro,noload,nodev,nosuid,noexec};
\item read private keys only through their role-specific configured paths; and
\item write signing requests, signatures, completed artifacts, logs, and temporary
files somewhere other than the signing medium.
\end{itemize}
If the filesystem is dirty, damaged, unexpectedly writable, or has the wrong
identity, stop signing. Do not replay its journal or repair it implicitly in
the routine path.
Permit writes to the decrypted volume, including a read-write mount, only
during an explicit signing-media maintenance operation, such as:
\begin{itemize}
\item initial key creation;
\item an authorized key or certificate rotation;
\item a verified metadata change;
\item primary-to-backup synchronization; or
\item filesystem repair.
\end{itemize}
A writable operation must identify whether it targets the primary or backup,
validate the pre-change key fingerprints and manifest, produce and validate a
post-change manifest, flush writes, cleanly unmount ext4, close LUKS2, and
detach the medium before it is complete.
Keep the backup detached except during planned synchronization, read-only
recovery testing, or repair. Routine trial-signature tests use the same
read-only policy as the primary.
The exact mount point, filesystem-cleanliness command, writable-maintenance
authorization, staging location, manifest format, synchronization procedure,
and temporary-data cleanup remain separate decisions.

\adrentry{0039}{Unlock both signing media with one shared passphrase}{Deferred for post-proof revalidation by ADR 0043}{2026-07-24}
Generate one high-entropy random passphrase and enroll it in a LUKS2 keyslot on
both signing-media volumes.
\begin{itemize}
\item Generate at least 128 bits of entropy with a cryptographically secure random
source and encode it in a form suitable for reliable manual entry or transfer.
\item Use the same passphrase for the primary and backup, while retaining the
independently generated LUKS2 volume master key on each medium.
\item Enter the passphrase interactively in the designated signing environment.
\item Do not store it on either signing medium, any Controller or Workstation, the
bootstrap or signing host, an installer, PXE storage, Git, a command line,
environment variable, shell history, log, or generated artifact.
\item Keep its authoritative custody off-host and physically separate from both
signing media.
\end{itemize}
Changing the shared passphrase is incomplete until the new credential has been
enrolled and tested on both media and the old credential has been removed from
both. Test each medium separately.
ADR 0040 subsequently selects the operator's existing off-Controller password
manager as the authoritative store. The passphrase generator and encoding,
human access policy, input method, keyslot layout, LUKS2 PBKDF parameters, and
rotation cadence remain separate decisions. ADR 0041 subsequently requires one
independent emergency copy, and ADR 0042 selects its sealed paper form.

\adrentry{0040}{Store the signing-media passphrase in an off-Controller password manager}{Deferred for post-proof revalidation by ADR 0043}{2026-07-24}
Use the operator's existing off-Controller password manager as the
authoritative store for the shared signing-media passphrase.
The password manager used for this purpose must:
\begin{itemize}
\item remain accessible when the Controller is powered off, damaged, isolated, or
being rebuilt;
\item not depend exclusively on DNS, identity, networking, or storage supplied by
the Controller;
\item protect its local or remote vault independently of the signing media; and
\item have an operator-tested recovery path independent of the Controller.
\end{itemize}
Retrieve the passphrase manually for a signing session and enter it
interactively into the designated signing environment. Do not add password
manager API credentials, automated vault retrieval, or a machine account to
the initial signing workflow. Do not cache or persist the retrieved passphrase
on the signing host.
The repository and non-secret machine records may contain only a stable,
non-secret reference describing which password-manager entry the operator must
retrieve. They must not contain the passphrase, a vault export, session token,
recovery code, or password-manager credential.
The specific password-manager product and account remain operator-owned
instance details rather than Controller profile dependencies. This design does
not require external registration; a qualifying local or externally hosted
manager is acceptable.
ADR 0041 subsequently requires one independent emergency copy, and ADR 0042
selects a sealed paper recovery sheet in locked, fire-resistant storage.
Password-manager entry naming, clipboard versus direct-entry handling,
authorized human access, audit cadence, and coordinated passphrase-rotation
procedure remain separate decisions.

\adrentry{0041}{Keep one independent emergency copy of the signing-media passphrase}{Deferred for post-proof revalidation by ADR 0043}{2026-07-24}
Maintain exactly one authorized emergency copy of the shared signing-media
passphrase outside the password manager.
The emergency copy must:
\begin{itemize}
\item reproduce the exact passphrase without depending on memory;
\item remain offline and require no Controller, home-network, DNS, identity, cloud,
or password-manager service;
\item be stored physically apart from the password manager's recovery material,
the primary signing medium, and the backup signing medium;
\item never appear in Git, ordinary documentation, an inventory database, a
Controller or Workstation, PXE storage, an installer, or a signing host;
\item be accessed only for password-manager recovery or a scheduled break-glass
test; and
\item be returned to protected storage or securely replaced immediately after use.
\end{itemize}
The password manager remains the normal authoritative source. The emergency
copy is a break-glass recovery path, not a second routine retrieval location.
A passphrase rotation is incomplete until the password manager, both LUKS2
volumes, and the emergency copy agree on the new credential and each has been
tested through its appropriate path. Destroy the obsolete emergency copy only
after the new one has been verified.
ADR 0042 subsequently selects a human-readable paper recovery sheet in a
tamper-evident sealed envelope kept in locked, fire-resistant storage. The
exact storage location, access control, test cadence, and destruction method
remain separate decisions.

\adrentry{0042}{Use a sealed paper emergency passphrase copy}{Deferred for post-proof revalidation by ADR 0043}{2026-07-24}
Print the exact signing-media passphrase as a human-readable recovery sheet.
Place it in a tamper-evident sealed envelope and keep the envelope in locked,
fire-resistant storage physically separate from:
\begin{itemize}
\item the primary signing medium;
\item the backup signing medium; and
\item the password manager's own recovery material.
\end{itemize}
The sheet may contain only:
\begin{itemize}
\item the exact passphrase in a transcription-resistant layout;
\item a non-secret purpose and version label;
\item enough non-secret instructions to identify the two signing-media LUKS2
volumes it unlocks; and
\item a verification value that detects transcription error without containing
another secret.
\end{itemize}
Do not print either signing private key, a password-manager credential, a vault
recovery code, or unrelated recovery material on the sheet.
Before sealing a new or rotated sheet, use it to unlock and close each medium
separately, verify both key fingerprints through a trial signature, and confirm
that no digital copy remains in the print path, spool, temporary directory, or
device storage under operator control.
Record only a non-secret envelope identifier, creation date, and seal-inspection
status outside the envelope. Do not record the passphrase or an image of the
sheet.
Opening the envelope is a break-glass event. Record the event without the
secret, replace the sheet and envelope after use, and rotate the passphrase if
custody or observation during the event is uncertain.
This plaintext physical sheet is an explicit exception to digital at-rest
encryption. Its locked location, separation, tamper evidence, and access
control form its protection boundary.
The exact storage location, authorized people, paper and printing method,
envelope type, verification-value format, inspection and test cadence, and
secure destruction method remain separate decisions.

\adrentry{0043}{Prove the functional environment before boot hardening}{Accepted}{2026-07-24}
Split delivery into two explicit milestones.
\#\#\# Milestone A: functional proof
Build and test the non-production provisioning path first:
\begin{itemize}
\item network boot into the interactive installer;
\item collect, validate, summarize, and explicitly confirm all destructive inputs;
\item record UEFI mode, Secure Boot capability and state, and TPM 2.0 availability
without changing firmware trust or enrolling a TPM token;
\item install the generic Controller profile with Arch, LUKS2, Btrfs, systemd-boot,
the two normal kernels, and the accepted shared-ESP layout;
\item exercise the permanent Windows maintenance installation and its independent
UEFI entry far enough to validate the functional dual-boot layout;
\item power off after installation, relocate to the isolated switch, and perform
the first installed boot;
\item after manual LUKS2 unlock, start Controller DHCP and DNS automatically;
\item verify isolated-network leases, \texttt{home.arpa} resolution, PXE service, and the
no-routing/no-NAT boundary;
\item exercise the Workstation profile's Controller-FQDN pivot; and
\item prove checkpoint, rollback, rebuild, and health-validation concepts without
depending on an offline signer.
\end{itemize}
Continue to verify the provenance and published checksums of downloaded
installer and package artifacts even though Milestone A does not prove the
final authenticated-boot chain.
Milestone A does not create, enroll, or use a custom lab Secure Boot identity,
TPM policy identity, TPM unlock token, signing medium, or signing credential.
It does not clear factory firmware keys, enter firmware Setup Mode, alter
firmware \texttt{db}, or implement TPM automatic LUKS2 unlock.
If the Arch or network-boot path cannot use an already trusted upstream chain,
Milestone A may run with Secure Boot disabled only under an explicit
\texttt{development-proof} mode. That mode must:
\begin{itemize}
\item display and record that authenticated boot is not being tested;
\item require a separate explicit operator acknowledgement before destructive
authorization;
\item use a manually entered LUKS2 passphrase at every Arch boot;
\item contain no production secrets or irreplaceable data; and
\item never be reported as a production-ready or fully accepted profile.
\end{itemize}
\texttt{development-proof} is a temporary project state recorded in the installation
manifest, not a permanent profile option that can be selected to avoid
hardening.
Windows installation and direct UEFI boot may be tested in Milestone A, but
final BitLocker, Secure Boot, TPM, and recovery behavior belongs to Milestone B.
\#\#\# Milestone B: security hardening and final acceptance
After Milestone A passes, revisit ADRs 0029 through 0042 as a group. They are
design input, not an instruction to enroll TPM state or manufacture keys,
credentials, or signing media now. Simplify, change, or replace them based on
the proven update and provisioning workflow.
Milestone B must then:
\begin{itemize}
\item enforce ADR 0020's Secure Boot and encryption requirements;
\item establish and test the selected network-boot, installed-boot, and recovery
trust chains;
\item implement and test the selected signing and recovery-key custody;
\item implement TPM automatic unlock only after its recovery path works;
\item complete Windows BitLocker and recovery acceptance; and
\item prove that an unattended cold restart restores Controller DHCP, DNS, and PXE
service; and
\item reprovision from the proven installer for final acceptance rather than
treating the development-proof installation as production.
\end{itemize}
No system passes final Controller or Workstation acceptance until Milestone B
is complete.

\adrentry{0044}{Use dnsmasq for initial Controller network services}{Accepted}{2026-07-24}
Use the Arch \texttt{dnsmasq} package as the initial host-native implementation of the
Controller's bundled network-services option.
When Controller network services are enabled:
\begin{itemize}
\item dnsmasq provides DHCPv4 and is the only address-assigning DHCP authority on
the managed layer-2 network;
\item dnsmasq serves the Controller's local \texttt{home.arpa} DNS data and eligible DHCP
lease or reservation names;
\item DHCP advertises the Controller as the DNS server and does not advertise a
default router on the isolated network;
\item dnsmasq supplies architecture-appropriate PXE information and a restricted,
read-only TFTP service for the small first-stage boot program; and
\item the service binds only to the explicitly selected managed interface and
address.
\end{itemize}
When external infrastructure owns DHCP and DNS, keep the installed Controller's
dnsmasq network-service instance stopped. Whether the temporary bootstrap host
uses UniFi boot options or a separate dnsmasq ProxyDHCP configuration remains a
later decision.
Do not deploy Kea, BIND, Unbound, CoreDNS, or a separate TFTP daemon for the
initial Controller network-services bundle. Serve iPXE scripts, kernels,
initramfs images, WIMs, and other substantial artifacts over a separately
selected HTTP(S) service rather than TFTP.
Revisit the service split if the proven environment requires DHCP high
availability, several routed networks and relays, database- or API-driven IP
address management, independently operated DNS, or authoritative DNS features
beyond dnsmasq's practical scope.

\adrentry{0045}{Collect a minimal explicit Controller IPv4 network plan}{Accepted}{2026-07-24}
When the bundled Controller network-services option is enabled, collect these
four machine-readable installation inputs:
\begin{itemize}
\item \texttt{managed\_ipv4\_cidr}: the canonical network address and prefix length for the
managed subnet;
\item \texttt{controller\_ipv4\_address}: the Controller's static service address;
\item \texttt{dhcp\_pool\_start}: the first address in the inclusive dynamic pool; and
\item \texttt{dhcp\_pool\_end}: the last address in the inclusive dynamic pool.
\end{itemize}
Derive rather than prompt for:
\begin{itemize}
\item the IPv4 netmask;
\item the network and broadcast addresses;
\item the client DNS-server address, which is
\texttt{controller\_ipv4\_address};
\item the DNS suffix, which is \texttt{home.arpa}; and
\item the absence of a default-router option.
\end{itemize}
Before destructive authorization, require all of the following:
\begin{itemize}
\item every input parses unambiguously as IPv4;
\item \texttt{managed\_ipv4\_cidr} uses the actual network address rather than an address
with host bits set;
\item the subnet contains enough usable addresses for the Controller and pool;
\item the Controller and both pool endpoints are usable unicast addresses inside
the subnet;
\item the pool start is not greater than the pool end;
\item the Controller address is outside the inclusive DHCP pool; and
\item the complete entered and derived network plan is displayed in the preflight
summary.
\end{itemize}
Record the confirmed inputs and derived plan in the non-secret installation
manifest. If Controller-owned DHCP and DNS are disabled, these managed-network
inputs are not required; the external-network host-addressing policy remains a
separate decision.
Do not prompt for a netmask, broadcast address, DNS server, DNS suffix, or
gateway separately.

\adrentry{0046}{Keep the homelab record inside Telos and split public profile from private instance}{Accepted}{2026-07-25}
Move the homelab design record, implementation and documentation into the Telos
repository, and split it in two:
\begin{itemize}
\item \textbf{Profile material is public.} Decision records, the generic Controller,
Workstation and Services profiles, the PXE design, the installer, the Ansible
roles, and both reconstruction manuals contain no instance data and publish to
the Telos site.
\item \textbf{Instance material is private.} Real hostnames, addresses, interface names,
disk serials, DHCP reservations, share paths and per-machine inventory live in
\texttt{homelab/instance/}, which is gitignored and never rendered into the site.
\end{itemize}
Documents that need an instance value read it from the private overlay at build
time and fall back to a clearly marked placeholder when the overlay is absent.
A published document must be complete and useful without the overlay.

\adrentry{0047}{Remove the permanent Windows maintenance installation from the Controller}{Accepted}{2026-07-25}
Remove the permanent Windows installation from the Controller profile.
\begin{itemize}
\item Perform firmware maintenance with \texttt{fwupd} where the vendor publishes to LVFS.
\item Otherwise use vendor-supplied bootable media, created on demand.
\item The Controller disk carries one ESP, one LUKS2 container and nothing else.
\item Preflight records firmware vendor, version and LVFS availability in the
installation manifest so the maintenance path is known before it is needed.
\end{itemize}
This decision governs the Controller profile only. The Workstation profile may
still install Windows; that is decided separately.

\adrentry{0048}{Serve boot and installation artifacts over HTTP with nginx}{Accepted}{2026-07-25}
Use the Arch \texttt{nginx} package as the Controller's artifact service.
\begin{itemize}
\item Serve a single read-only document root, \texttt{/srv/http/boot}.
\item Bind only to the managed interface and the Controller service address.
\item Serve plain HTTP on the isolated proof network. iPXE's HTTPS support depends
on build options and an embedded trust store, and the proof network has no
certificate authority; artifact integrity comes from published checksums
verified by the installer, not from transport.
\item Publish a \texttt{manifest.json} beside the artifacts listing every file with its
SHA-256. The installer verifies each artifact against it before use.
\item Run as a distinct system user with no write access to the document root.
\end{itemize}
Do not use a general-purpose application server, and do not enable directory
autoindex on the published root.

\adrentry{0049}{Build the installer as an Archiso netboot image driving a Python installer}{Accepted}{2026-07-25}
Build the provisioning environment as a custom Archiso profile, booted over the
network, running a Python installer from this repository.
\begin{itemize}
\item The image carries the tools the installer needs and nothing else: \texttt{parted},
\texttt{cryptsetup}, \texttt{btrfs-progs}, \texttt{dosfstools}, \texttt{arch-install-scripts},
\texttt{systemd-ukify}, \texttt{python}, \texttt{nginx}-less, no desktop.
\item The installer is ordinary Python 3 using only the standard library, so it can
be unit-tested off-target. \texttt{homelab/lib/} holds the logic; \texttt{homelab/bin/}
holds the entry points.
\item Only administrator \textbf{public} SSH keys are baked into the image, enabling a
parallel SSH session for observation. No private key, token or credential is
ever placed in an image.
\item The image is built reproducibly from a pinned package list and a recorded
mirror snapshot, and its SHA-256 is published in the artifact manifest.
\end{itemize}
Do not use WDS, MDT, or a vendor deployment suite.

\adrentry{0050}{Select the managed interface explicitly and pin it by MAC address}{Accepted}{2026-07-25}
Select the managed interface as an explicit installation input.
\begin{itemize}
\item Preflight enumerates every physical Ethernet interface with its kernel name,
permanent MAC address, current link state and observed speed.
\item The operator selects one. There is no default and no automatic choice.
\item The installer records the \textbf{permanent MAC address} as the identity and writes
a systemd \texttt{.link} file pinning a stable name, \texttt{lan0}, to that MAC.
\item dnsmasq, nginx and the static address configuration all reference \texttt{lan0}.
\item First-boot activation verifies that \texttt{lan0} exists, carries the expected MAC and
has carrier before starting network services, and fails closed otherwise.
\item Wireless interfaces are not offered.
\end{itemize}

\adrentry{0051}{Fix the Controller disk layout and ESP size}{Accepted}{2026-07-25}
Use this GPT layout on a single Controller disk:
\begin{tabularx}{\linewidth}{@{}lllL@{}}
\toprule
\textbf{\#} & \textbf{Size} & \textbf{Type} & \textbf{Contents}\\
\midrule
1 & 2 GiB & EFI System (FAT32) & systemd-boot, three UKIs, update headroom\\
2 & rest & Linux LUKS2 & Btrfs, subvolumes per ADR 0027\\
\bottomrule
\end{tabularx}
\begin{itemize}
\item 2 GiB is deliberate headroom, not a measurement. It costs nothing on any
plausible Controller disk and removes an entire class of failed-update
outage. The ESP is the one partition that cannot be grown after the fact
without moving everything after it.
\item No XBOOTLDR, no Microsoft reserved partition, no Windows recovery partition.
\item Swap is a file inside the encrypted Btrfs root, not a separate partition, so
ADR 0020's encrypted-swap requirement is satisfied without a second LUKS
container.
\item Hibernation is disabled on the Controller. A machine whose purpose is to
answer DHCP does not suspend.
\item The installer refuses any disk smaller than 64 GiB and requires the operator
to confirm the disk by model, serial and size, never by \texttt{/dev/sdX}.
\end{itemize}

\adrentry{0052}{Build the first Controller from an existing workstation}{Accepted}{2026-07-25}
Use an existing Arch workstation as the temporary bootstrap host for the first
Controller build.
\begin{itemize}
\item Run the same \texttt{dnsmasq} and \texttt{nginx} configurations the Controller profile uses,
generated by the same code, from a clearly labelled bootstrap directory.
\item On a network that already has a DHCP authority, run dnsmasq in \textbf{ProxyDHCP}
mode so it supplies only boot options and never addresses. Two address-
assigning authorities on one layer-2 network remain forbidden by ADR 0008.
\item The bootstrap host builds the Archiso image and publishes the artifact
manifest.
\item Bootstrap state is temporary. After the Controller passes acceptance, the
bootstrap services are stopped and the Controller serves subsequent installs.
\end{itemize}
The bootstrap host must not later be provisioned as a managed machine while it
is the only source of the boot chain.

\adrentry{0053}{Own day-two configuration with Ansible from this repository}{Accepted}{2026-07-25}
Split the lifecycle in two.
\textbf{Install time} does only what cannot be done later: partition, encrypt, create
the filesystem, install the base system and kernels, write the boot artifacts,
configure the managed interface and static address, install \texttt{sshd} with the
administrator public key, and record the manifest. Nothing else.
\textbf{Converge time} is Ansible, run over SSH from this repository, and owns
everything else: package sets, dnsmasq and nginx configuration, the services
role, identity client configuration, users, sudo policy, health checks and
update orchestration.
\begin{itemize}
\item Inventory lives in the private overlay; roles and playbooks are public.
\item Every role must be idempotent and must pass \texttt{---check} cleanly against a
converged host.
\item No secret is stored in a playbook. Secret material is referenced, and the
mechanism is chosen in a later ADR.
\item The Controller converges itself the same way a Workstation does. There is no
privileged manual path.
\end{itemize}

\adrentry{0054}{Define an optional Services role for always-on applications}{Accepted}{2026-07-25}
Define a third role, \texttt{services}, separate from Controller and Workstation.
\begin{itemize}
\item Implement each application as a Podman \textbf{quadlet} --- a systemd unit
generated from a declarative container definition --- so lifecycle, ordering,
restart policy and journal integration are ordinary systemd.
\item Applications that need discovery run with host networking, which is required
for HomeKit and UPnP to work at all.
\item USB radios are bound by stable \texttt{by-id} path and declared per application.
\item Application state lives on its own Btrfs subvolume, excluded from the
operating-system checkpoint under ADR 0027 and backed up separately, because
its restore point is not the same as the operating system's.
\item The role is \textbf{disabled by default} and enabled per instance.
\item The role may be assigned to the Controller host today. Moving it to its own
machine later changes which host claims the role and nothing else.
\end{itemize}
The Controller profile does not depend on the services role, and the services
role does not depend on Controller-hosted DHCP or DNS beyond ordinary name
resolution.

\adrentry{0055}{Use Samba AD DC for shared Windows and Linux identity}{Accepted}{2026-07-25}
Use Samba Active Directory Domain Services as the directory.
\begin{itemize}
\item Run it in a dedicated VM, not on the Controller host. It is a distinct trust
boundary, it needs its own restart and upgrade cadence, and ADR 0014 permits a
VM exactly when those conditions hold.
\item Arch clients use SSSD's AD provider with credential caching.
\item Windows Pro clients domain-join. Windows Home cannot and keeps local accounts.
\item Home directories stay local; \texttt{pam\_mkhomedir} creates them on first login.
\item Every managed machine keeps a separately named local break-glass administrator
with its own sudo rule, and that account is never a directory account.
\item UID and GID come from the directory. The existing local \texttt{ksh} at UID 1000 is
migrated deliberately and only after piloting with a separate test identity.
\item Set an explicit offline credential lifetime longer than a normal trip rather
than relying on a default.
\end{itemize}
A second domain controller on separate physical hardware is required before the
directory is considered production, and is deferred until the first one works.

\adrentry{0056}{Prove every change in a QEMU/OVMF matrix before hardware}{Accepted}{2026-07-25}
Maintain a QEMU/OVMF acceptance matrix in \texttt{homelab/qemu/}, runnable with one
command.
The matrix builds a virtual isolated network with no router and boots:
\begin{itemize}
\item a virtual \textbf{bootstrap} host serving ProxyDHCP, TFTP and HTTP;
\item a virtual \textbf{Controller} target, installed from the real installer with no
human present; ADR 0058 defines the mechanism, which drives the genuine
interactive prompts through a pseudo-terminal rather than adding an
unattended code path; and
\item a virtual \textbf{Workstation} target that pivots on the Controller's FQDN.
\end{itemize}
It then asserts, automatically:
\begin{itemize}
\item the installer refuses every invalid network plan in its rejection corpus;
\item the Controller installs, powers off, and activates services on first boot;
\item a client receives a lease inside the configured pool and never the Controller
address;
\item \texttt{home.arpa} resolves the Controller;
\item \textbf{no default route is advertised}, per ADR 0011;
\item the artifact manifest checksums verify;
\item Ansible converge is idempotent --- a second run reports no changes; and
\item a checkpoint can be created, rolled back to, and validated.
\end{itemize}
Physical installation is permitted only after the matrix passes. UEFI variables
use a per-run OVMF variable store so firmware state never leaks between runs.

\adrentry{0057}{Record corrections to the deferred signing design}{Accepted}{2026-07-25}
Do not act on this ADR during Milestone A. When ADRs 0029 through 0042 are
revalidated, treat the following as required inputs.
\textbf{1. The key separation in ADRs 0033 to 0035 is largely notional.} ADR 0033
makes the TPM policy key cryptographically distinct from the Secure Boot key,
which is correct. ADR 0035 then stores both on the same removable medium and ADR
0039 gives the primary and backup media the same passphrase. One compromised
passphrase yields both identities on both media, so the separation survives on
paper only. Either separate the custody or stop paying for the separation.
\textbf{2. ADR 0039's shared passphrase defeats the backup's independence.} A backup
exists to survive a failure the primary did not. Sharing its passphrase means a
credential compromise takes both copies simultaneously.
\textbf{3. ADR 0030's additive trust model may be infeasible on the actual hardware.}
Many consumer and OEM firmwares permit \texttt{db} modification only after entering
Setup Mode, which clears PK and therefore replaces platform ownership --- the
exact thing ADR 0030 forbids. ADR 0030 handles this by failing preflight, which
means the profile may simply not support the machine. Before Milestone B,
test enrollment on the real target firmware and decide what happens when it
cannot be done additively.
Also revisit whether three signed UKIs and two signing identities are
proportionate to a single-operator homelab, given the operational cost the
custody chain in ADRs 0036 through 0042 imposes.

\adrentry{0058}{Drive acceptance tests through the real interactive installer}{Accepted}{2026-07-26}
\textbf{No unattended installation code path exists.} The installer is interactive
and only interactive.
The acceptance matrix drives it by attaching to a pseudo-terminal and answering
prompts the way a person would, including the final confirmation.
\begin{itemize}
\item Prompt definitions live in one registry, \texttt{homelab/lib/prompts.py}: each prompt
has a stable identifier, its human-facing text, its validator and its help
text.
\item The installer renders that registry. The test harness imports the same
registry to know what to expect. There is one source of truth, so changing a
prompt's wording cannot silently desynchronize the tests.
\item The harness matches on prompt identity, supplies an answer, and asserts on
what the installer prints back.
\item There is no flag, environment variable or file that skips a prompt, and no
code path that reaches a destructive operation without the confirmation
having been answered.
\end{itemize}
The final confirmation requires the operator to **type the target disk's serial
number**, not to answer yes. A serial cannot be typed from muscle memory, cannot
be answered correctly by someone who has not read the summary, and cannot be
satisfied by a stray keystroke.

\adrentry{0059}{Fail fast, leave the evidence, do not roll back or resume}{Accepted}{2026-07-26}
\textbf{Stop at the failing step. Change nothing further. Report precisely.}
\begin{itemize}
\item No rollback. A cleanup path is a second destructive code path that only ever
executes when something is already wrong, and running it destroys the state
that explains the failure.
\item No resume. Resume logic must reason correctly about state it did not observe,
which is where installers grow their worst bugs. A re-run starts from the
beginning.
\item On failure the installer prints: the steps that completed, the step that
failed, the exact command, its exit status, and its captured output; then
states plainly that the disk is \textbf{not bootable} and that the fix is to run
the installer again.
\item The installer exits non-zero so the acceptance harness can detect it.
\end{itemize}
Steps are declared as either non-destructive or destructive. **A destructive
step cannot execute without an authorization token**, and that token can only be
produced by the confirmation described in ADR 0058. This is structural, not a
convention: there is no boolean to pass by mistake.

\adrentry{0060}{Record a non-secret manifest on the installed system}{Accepted}{2026-07-26}
Write one JSON manifest to \texttt{/etc/homelab/manifest.json} on the installed root,
and echo it to the console at the end of the run.
\begin{itemize}
\item The machine carries its own provenance. A Controller found in a cupboard in
three years can state what it is, when it was built, and from what.
\item The console copy is how the acceptance harness under ADR 0056 captures the
manifest without mounting an encrypted filesystem.
\item Nothing is uploaded during installation. There is no service to upload to when
the first Controller is built, and adding one would make provisioning depend
on it.
\item The manifest is \textbf{non-secret by construction} and is validated against that
before being written. It records identities, not credentials: no passphrase,
no recovery key, no private key, no token, no LUKS header, no password hash.
\item It is not written to the ESP. ADR 0020 keeps the ESP minimal and secret-free,
and interface MACs, disk serials and network plans are exactly the instance
data ADR 0046 keeps out of anything published or casually readable.
\end{itemize}
Recorded: schema version; installer version and the artifact checksums it
verified; timestamp; profile; hostname; whether this was a \texttt{development-proof}
installation; firmware mode, Secure Boot state and TPM presence as *observed*;
target disk path, model, serial and size; managed interface name and permanent
MAC; the confirmed network inputs and every derived value; and the partition
layout actually written.

\adrentry{0061}{Test the disk-writing steps against loopback devices}{Superseded by ADR 0062}{2026-07-26}
Write the disk steps so they take a \textbf{device path} and nothing else about how
that device came to exist. A test can then supply a loopback device backed by a
sparse file, and the genuine \texttt{sgdisk}, \texttt{cryptsetup} and \texttt{mkfs.btrfs} run against
it for real.
\begin{itemize}
\item Loopback tests live in \texttt{homelab/tests/integration/} and are \textbf{not} part of
the default suite. They need root and they are slow.
\item Run them with an explicit target: \texttt{make homelab-integration}.
\item Each test creates its own sparse backing file in a temporary directory,
attaches it, runs the step, asserts on the result by re-reading the device
with \texttt{sgdisk ---print}, \texttt{cryptsetup luksDump} and \texttt{btrfs subvolume list}, and
detaches and deletes it unconditionally.
\item A test must never touch a device it did not create. The helper refuses any
path that is not a loop device it attached itself, and the refusal is tested.
\item \texttt{pacstrap} and UKI generation are \textbf{not} covered here. They need network
access and a full Arch environment, so they belong to the QEMU matrix under
ADR 0056.
\end{itemize}
This does not replace the acceptance matrix. Loopback tests answer "does this
step produce the layout it claims"; the matrix answers "does a machine built
this way boot and serve DHCP".

\adrentry{0062}{Reach a running install cycle before hardening the installer}{Accepted}{2026-07-26}
Reach a running install cycle first. Specifically:
\begin{itemize}
\item ADR 0061 is superseded. There is no loopback test tier.
\item \texttt{disks.py} provides the commands to run and nothing else. The \texttt{verify\_*}
functions that re-read and parse \texttt{sgdisk}, \texttt{cryptsetup} and \texttt{btrfs} output are
removed; in the QEMU matrix a wrong layout presents as a machine that does not
boot, within about ninety seconds.
\item Priority order is now: installer entry point, Archiso profile, QEMU matrix,
then a real install on the spare machine.
\item Hardening resumes after the cycle exists, directed by what actually fails.
\end{itemize}
Retained deliberately, because retrofitting them is expensive and they are what
stands between an operator error and an erased disk:
\begin{itemize}
\item the validation layer ---- \texttt{netplan}, \texttt{prompts}, \texttt{preflight};
\item the structural authorization token in \texttt{steps.py}; and
\item the manifest's non-secret guarantee.
\end{itemize}

\end{document}
